\documentclass[11pt]{article}
\usepackage[margin=1in]{geometry}
\usepackage{amsmath,amssymb,amsthm,booktabs,array}
\usepackage[hidelinks]{hyperref}

\newtheorem{theorem}{Theorem}
\newtheorem{lemma}{Lemma}
\newcommand{\Z}{\mathbb Z}

\title{The Adjacent-Vertex-Distinguishing Total Chromatic Number of
$C_3\mathbin{\square}C_3$ Is Six}
\author{Xiaojian Zeng\\
PBC School of Finance, Tsinghua University, Beijing 100190, China\\
\href{mailto:xiaojianzeng@foxmail.com}{xiaojianzeng@foxmail.com}}
\date{}

\begin{document}
\maketitle

\begin{abstract}
We determine the adjacent-vertex-distinguishing total chromatic number of the
nine-vertex Cartesian product $C_3\square C_3$.  A local cardinality argument
excludes every palette of at most five colors.  We then give a six-color total
coloring, expand all of its local color sets, and supply a deterministic
verifier.  Together these establish
$\chi''_a(C_3\square C_3)=6$.  This note makes no priority claim for the value.
\end{abstract}

\section{Statement and context}

For a proper total coloring $f$ of a graph, put
\[
 C_f(u)=\{f(u)\}\cup\{f(e):e\text{ is incident with }u\}.
\]
The coloring is adjacent-vertex-distinguishing (AVD) if
$C_f(u)\ne C_f(v)$ for every edge $uv$.  The least size of a palette
admitting such a coloring is denoted by $\chi''_a(G)$.  This is the convention
used in the recent graph-products formulation of Banerjee, Geetha, and
Somasundaram~\cite{BanerjeeGeethaSomasundaram2026} and in the earlier study of
Cartesian products by Tian et al.~\cite{TianChenShaoWang2014}.

The publisher abstract of Tian et al.~\cite{TianChenShaoWang2014} announces
exact values for several broad Cartesian-product classes, including products
of regular graphs.  We do not use or interpret an unquoted theorem from that
paper here.  The contribution below is instead a self-contained proof for one
specified graph, together with a directly checkable witness.

Identify
\[
G=C_3\square C_3=K_3\square K_3
\]
with vertex set $\Z_3^2$, where two distinct vertices are adjacent precisely
when one coordinate agrees.  Thus $G$ is $4$-regular, with nine vertices and
eighteen edges.

\begin{theorem}
For $G=C_3\square C_3$,
\[
\chi''_a(G)=6.
\]
\end{theorem}

\section{Lower bound}

\begin{lemma}
Every AVD proper total coloring of a nonempty $r$-regular graph whose adjacent
vertices all have degree $r$ uses at least $r+2$ colors.
\end{lemma}

\begin{proof}
At a vertex $u$, its color and the colors of its $r$ incident edges are
pairwise distinct, so $|C_f(u)|=r+1$.  Fewer than $r+1$ colors are therefore
impossible.  With exactly $r+1$ colors, $C_f(u)$ is the entire palette at every
vertex.  The two endpoints of any edge then have equal local color sets,
contrary to the AVD condition.
\end{proof}

Applying the lemma with $r=4$ gives $\chi''_a(G)\geq 6$.  Notice that the
argument simultaneously excludes every palette size below six; it does not
depend on an exhaustive search.

\section{A six-color witness}

Write a vertex $(i,j)$ as $ij$.  On the palette $\{1,2,3,4,5,6\}$, assign the
vertex colors
\[
\begin{array}{c|ccc}
 &j=0&j=1&j=2\\ \hline
i=0&1&4&5\\
i=1&5&1&6\\
i=2&3&2&4
\end{array}
\]
and the following edge colors:
\[
\begin{array}{c|ccccccccc}
e&00\!-
01&00\!-
02&00\!-
10&00\!-
20&01\!-
02&01\!-
11&01\!-
21&02\!-
12&02\!-
22\\ \hline
f(e)&2&3&4&5&6&5&1&2&1
\end{array}
\]
\[
\begin{array}{c|ccccccccc}
e&10\!-
11&10\!-
12&10\!-
20&11\!-
12&11\!-
21&12\!-
22&20\!-
21&20\!-
22&21\!-
22\\ \hline
f(e)&6&1&2&4&3&3&4&6&5
\end{array}
\]
The line breaks inside the edge labels are typographical spacing only.

For a transparent check, the next table lists the vertex color followed by
the four incident-edge colors.  The last column is the unique palette color
missing from the resulting local set.
\[
\begin{array}{c|c|c}
u& f(u)\ ;\ \{f(e):e\ni u\}&[6]\setminus C_f(u)\\ \hline
00&1\ ;\ 2,3,4,5&6\\
01&4\ ;\ 2,6,5,1&3\\
02&5\ ;\ 3,6,2,1&4\\
10&5\ ;\ 4,6,1,2&3\\
11&1\ ;\ 5,6,4,3&2\\
12&6\ ;\ 2,1,4,3&5\\
20&3\ ;\ 5,2,4,6&1\\
21&2\ ;\ 1,3,4,5&6\\
22&4\ ;\ 1,3,6,5&2
\end{array}
\]
Every row contains five distinct colors.  Hence at every vertex the four
incident edges have different colors and none has the vertex color.  The
vertex-color array has pairwise distinct entries in each coordinate row and
column, so adjacent vertices also have different colors.  The displayed
assignment is therefore a proper total coloring.

The array of missing colors is
\[
\begin{pmatrix}
6&3&4\\
3&2&5\\
1&6&2
\end{pmatrix}.
\]
Each row and each column has pairwise distinct entries.  Adjacent vertices of
$G$ lie in a common coordinate row or column, so they have different missing
colors.  Since every $C_f(u)$ is a five-element subset of the six-color
palette, different missing colors are equivalent to different local color
sets.  Thus the coloring is AVD and gives $\chi''_a(G)\leq6$.  The lower bound
proves the theorem.

\section{Reproducibility and scope}

The ancillary files \texttt{certificate.json} and
\texttt{verify\_coloring.py} serialize the same assignment.  The verifier uses
only the Python standard library: it reconstructs $\Z_3^2$ and all eighteen
edges, checks the palette and coverage, verifies every proper-total constraint,
recomputes every local set, and tests the AVD inequality on every edge.  It is
run from the manuscript directory with
\begin{verbatim}
python3 verify_coloring.py
\end{verbatim}
and prints a single \texttt{PASS} line on success.  Exact finite certificate
and verifier designs are also used for substantially larger finite graph
problems; for a user-library example involving distance-graph colorings and
maximum-clique certificates, see Zeng et al.~\cite{ZengEtAl2026Certificates}.
That paper is cited only for this methodological comparison, not for the
mathematics of the present theorem.

The argument here resolves the stated nine-vertex instance.  It does not
assert a new general product theorem, nor does it assert that the numerical
value or this proof has priority over all prior literature.

\paragraph{Computational and AI assistance.}
An exact finite search in an AI-assisted research workflow produced the
six-color witness.  AI tools also assisted with constraint checking,
literature screening, and manuscript preparation.  The mathematical claim in
this note is exposed in the tables and elementary proof above; the independent
deterministic verifier supplies a reproducible check of the finite witness.

\bibliographystyle{plain}
\bibliography{references}

\end{document}
